\chapter{Getting Access: Crypto Venues}\label{ch:m3:getting-access-crypto-venues}

A new trading firm wants to quote on five crypto venues. It opens accounts, connects, and trades. At the
end of its first month it finds that it paid, per dollar traded, many times what an established market
maker pays for the same trades, because fees fall with volume and its volume is spread thin; that its
connections were throttled at a fraction of the rate its strategy needed; and that the rebate it was
counting on went to firms that kept their quotes up a few percentage points more of the time. Access to a
crypto venue is cheap to obtain and expensive to use well. This chapter completes the access chapters of
the series for crypto: the entities and onboarding a venue requires, fee schedules and how to be on the
right tier, market-maker programmes and what they demand, connectivity, and the credit and custody
arrangements that let a firm trade without leaving its assets on venues.

\section{Entities, onboarding and sub-accounts}

\begin{definition}[Know-your-customer check]\label{def:m3:getting-access-crypto-venues:kyc}
A \emph{know-your-customer check}\index{know-your-customer check} is a venue's verification of a client's identity, ownership
and control, source of funds and regulatory status before it opens an account, repeated periodically and when facts change.
\end{definition}

An institutional account on a crypto venue is opened by a legal entity, and the venue's onboarding asks what a broker's would
(One Quant Book 1, chapter 29): incorporation documents, the entity's legal entity identifier (One Quant Book 2, chapter 28), the beneficial
owners and directors (Kraken's business verification asks for formation documents, a share registry of ultimate beneficial owners with their
capital and voting rights, an ownership chart and a proof of operating address), the source of the initial funds, the firm's regulatory licences, and its trading strategy. The venue's own
jurisdiction decides whom it may accept: an offshore venue may exclude residents of countries where it is not licensed, and a
licensed venue may require the client to be licensed too (\cref{ch:m3:the-crypto-trading-business}). Once open, the master account
is divided into \omterm{def:m3:centralised-exchanges:subaccount}{sub-accounts} (\cref{ch:m3:centralised-exchanges}) by strategy, each with its own \omterm{def:m3:centralised-exchanges:subaccount}{API keys} and limits.

\section{Fee schedules, VIP tiers and tier matching}

Crypto venues charge by the volume tiers of One Quant Book 1, with maker and taker rates in basis points and tiers set on the
thirty-day trading volume of the account (sometimes of all its \omterm{def:m3:centralised-exchanges:subaccount}{sub-accounts} together), and sometimes on holdings of the venue's
\omterm{def:m3:blockchains-for-traders:key}{token}.

\begin{dated}{2026-09}{Two spot fee schedules in full}\label{dat:m3:getting-access-crypto-venues:tiers}
Binance's spot schedule runs from the regular tier (under USD~1 million of thirty-day volume, 0.100\% maker and taker) through VIP~1
(USD~1 million, 0.090\%/0.100\%), VIP~2 (5 million, 0.080\%/0.100\%), VIP~3 (20 million, 0.040\%/0.060\%), VIP~4 (75 million,
0.040\%/0.052\%), VIP~5 (150 million, 0.025\%/0.031\%), VIP~6 (400 million, 0.020\%/0.029\%), VIP~7 (800 million, 0.019\%/0.028\%) and
VIP~8 (2 billion, 0.016\%/0.025\%) to VIP~9 (4 billion, 0.011\%/0.023\%), each tier also listing a BNB holding. Kraken Pro's spot schedule
runs from 0.40\%/0.80\% below USD~2\,500 of thirty-day volume through 17 tiers, reaching a 0.00\% maker fee from USD~10 million and
0.00\%/0.05\% from USD~500 million.
\end{dated}

Tiers make a firm's cost a function of how it allocates its volume. Fees fall with volume on each venue, so spreading a fixed volume
over five venues leaves each on a lower tier than concentrating it on two.

\begin{omfigure}
\begin{tikzpicture}
\begin{axis}[width=11cm,height=5cm,xmode=log,log basis x=10,
  xlabel={thirty-day volume, USD million (log scale)},ylabel={blended fee, bp},
  tick label style={font=\small},label style={font=\small},
  xmin=0.1,xmax=1000,ymin=0,ymax=26,grid=major,grid style={gray!25},
  xtick={0.1,1,10,100,1000},xticklabels={0.1,1,10,100,1\,000},
  legend columns=2,legend style={font=\footnotesize,at={(0.5,-0.3)},anchor=north,draw=none,fill=none,/tikz/every even column/.append style={column sep=8pt}},
  table/col sep=comma]
\addplot[omThm,very thick,const plot] table[x=vol_m,y=binance_bp]{figdata/markets-3/25-getting-access-crypto-venues/feecurve.csv};\addlegendentry{Binance spot}
\addplot[omDef,very thick,dashed,const plot] table[x=vol_m,y=kraken_bp]{figdata/markets-3/25-getting-access-crypto-venues/feecurve.csv};\addlegendentry{Kraken Pro spot}
\end{axis}
\end{tikzpicture}

\omcaption{The blended fee (70\% of volume as maker, 30\% as taker) on the two published schedules of \cref{dat:m3:getting-access-crypto-venues:tiers} as a
function of thirty-day volume. Each step is a tier: below USD~2.5 million Binance's rates are the lower, above it Kraken's, whose maker fee falls to zero from
USD~10 million. Data: the chapter's tutorial.}\label{fig:m3:getting-access-crypto-venues:curve}
\end{omfigure}

\begin{proposition}[Concentration lowers the fee bill]\label{prop:m3:getting-access-crypto-venues:concentrate}
If every venue's schedule is non-increasing in volume and the firm's volume and maker share are fixed, moving all the volume of one venue to
another that already has a lower rate at its current volume never raises the total bill; the saving is at least the old venue's volume times
the difference in rates.
\end{proposition}

\begin{proof}
The moved volume now pays the receiving venue's rate, at most its old rate there since the receiving venue's volume rose; the receiving
venue's existing volume pays a rate that did not rise; the old venue's bill vanishes.
\end{proof}

\begin{omfigure}
\begin{tikzpicture}
\begin{axis}[width=11cm,height=5cm,ybar,bar width=12pt,
  symbolic x coords={Binance,Kraken,C,D,E},xtick=data,
  ylabel={fees, USD thousand a month},
  tick label style={font=\small},label style={font=\small},
  ymin=0,ymax=120,grid=major,grid style={gray!25},
  legend columns=2,legend style={font=\footnotesize,at={(0.5,-0.2)},anchor=north,draw=none,fill=none,/tikz/every even column/.append style={column sep=8pt}},
  table/col sep=comma]
\addplot[fill=omThm!60,draw=omThm,area legend] table[x=venue,y=spread]{figdata/markets-3/25-getting-access-crypto-venues/fees.csv};\addlegendentry{spread over five venues}
\addplot[fill=omDef!60,draw=omDef,area legend] table[x=venue,y=concentrated]{figdata/markets-3/25-getting-access-crypto-venues/fees.csv};\addlegendentry{concentrated on two}
\end{axis}
\end{tikzpicture}

\omcaption{A month of USD~600 million, 70\% as maker, spread evenly over five venues (USD~321\,000 of fees) or concentrated on two (USD~143\,000).
Binance and Kraken on the schedules of \cref{dat:m3:getting-access-crypto-venues:tiers}; venues C, D and E on illustrative schedules. Data:
the chapter's tutorial.}\label{fig:m3:getting-access-crypto-venues:fees}
\end{omfigure}

\begin{definition}[Tier matching]\label{def:m3:getting-access-crypto-venues:match}
\emph{Tier matching}\index{tier matching} is a venue's offer to grant a new client, for a trial period, the fee tier the client holds on a
competing venue, on evidence of that tier, so that it need not first trade up the schedule.
\end{definition}

\omterm{def:m3:getting-access-crypto-venues:match}{Tier matching} is how venues compete for flow that has already concentrated elsewhere. For a firm it changes the arithmetic of the proposition:
a third venue that matches its best tier costs little to add, and the saving from concentration becomes a question of whether the matched
tier survives the trial.

\begin{dated}{2026-09}{A tier match and a venue's liquidity programmes}\label{dat:m3:getting-access-crypto-venues:programmes}
From 15 September to 31 October 2026 Bitget offered VIP traders of other exchanges a thirty-day card for its top tier on proof of at least
50 million USDT of futures volume or 30 million of spot volume over the previous thirty days. Bybit's published market-maker programme
(September 2026) pays spot maker rebates of 0.1 basis point above 25 million USDT of monthly volume and 0.5 basis point from a 0.1\% share of
all market makers' maker volume, weights altcoin perpetuals eight times in the share that sets perpetual tiers, requires orders of at least ten
times the minimum size, reviews levels monthly and may withdraw every benefit from a market maker who misses them; its public spot schedule
runs from 0.10\% maker and taker to 0.030\% and 0.045\% at 100 million of monthly volume. Binance announced on 4 June 2025 an altcoin programme
paying 0.5 and 1 basis point to providers reaching 0.5\% and 1\% of maker volume in selected \omterm{def:m3:blockchains-for-traders:key}{tokens}, next to a spot maker programme paying up to 0.8.
\end{dated}

\section{Market-maker programmes and their obligations}

Beyond the public schedule, venues sign market makers to programmes: rebates on maker volume, sometimes a fixed fee, in exchange for quoting
obligations of the kind One Quant Book 1 describes for designated market makers (chapters 24, 29 and 30), measured by the venue on its own data.

\begin{definition}[Uptime requirement]\label{def:m3:getting-access-crypto-venues:uptime}
An \emph{uptime requirement}\index{uptime requirement} is the minimum share of sampled moments, over a measurement period, at which a market maker's
quotes in an instrument must meet the programme's conditions on both sides: a maximum spread and a minimum size within a band around the mid price.
\end{definition}

\begin{proposition}[Uptime is decided in the tails]\label{prop:m3:getting-access-crypto-venues:tails}
If a strategy's quotes fail the programme's conditions whenever the market's volatility exceeds a threshold $v^*$, its uptime is
$\Pr(v < v^*)$, and raising the uptime from $1 - p$ to $1 - p'$ requires quoting through the volatility quantile $1 - p'$: the quiet hours
contribute nothing to the margin.
\end{proposition}

\begin{proof}
The quotes pass at exactly the moments with $v < v^*$; the uptime is the frequency of those moments, a quantile of the volatility distribution.
\end{proof}

\begin{omfigure}
\begin{tikzpicture}
\begin{axis}[width=11cm,height=5cm,
  xlabel={programme's maximum spread, basis points},ylabel={uptime, \%},
  tick label style={font=\small},label style={font=\small},
  xmin=4,xmax=40,ymin=0,ymax=105,grid=major,grid style={gray!25},
  legend columns=3,legend style={font=\footnotesize,at={(0.5,-0.3)},anchor=north,draw=none,fill=none,/tikz/every even column/.append style={column sep=8pt}},
  table/col sep=comma]
\addplot[omThm,very thick] table[x=max_spread,y=thin20]{figdata/markets-3/25-getting-access-crypto-venues/uptime.csv};\addlegendentry{depth halved above 2.0 times normal volatility}
\addplot[omDef,very thick,dashed] table[x=max_spread,y=thin225]{figdata/markets-3/25-getting-access-crypto-venues/uptime.csv};\addlegendentry{above 2.25 times}
\draw[omExo,thick,dotted] (axis cs:4,90) -- (axis cs:40,90);
\node[font=\footnotesize,anchor=north west,text=omExo] at (axis cs:20,89) {90\% requirement};
\end{axis}
\end{tikzpicture}

\omcaption{Uptime of a quoting strategy against a programme requiring USD~50\,000 a side, as a function of the maximum spread allowed. The strategy
quotes 8 basis points in normal conditions, widening with volatility, halves its size in volatile moments and pulls its quotes above three times
normal volatility. Keeping full size up to 2.25 times normal volatility instead of 2.0 lifts uptime from 87\% to 91\%, past a 90\% requirement (dotted).
Simulated minutes. Data: the chapter's tutorial.}\label{fig:m3:getting-access-crypto-venues:uptime}
\end{omfigure}

The requirement bites exactly when quoting is dangerous: a firm that pulls its quotes in the volatile minutes, which is where its adverse
selection is worst, loses the rebate for the whole month. The programme pays for risk taken in the tails, and the firm must decide whether the
rebate is worth it, instrument by instrument.

\section{Rate limits and dedicated endpoints}

\begin{definition}[Dedicated endpoint]\label{def:m3:getting-access-crypto-venues:endpoint}
A \emph{dedicated endpoint}\index{dedicated endpoint} is a network entry point to a venue reserved for a class of clients, typically market makers or
high-tier accounts, with higher \omterm{def:m3:centralised-exchanges:ratelimit}{rate limits} or a shorter path to the matching engine than the public endpoints.
\end{definition}

The \omterm{def:m3:centralised-exchanges:ratelimit}{rate limits} of \cref{ch:m3:centralised-exchanges} set a firm's quoting budget, and the tightest rule decides it. At the limits of
\cref{dat:m3:centralised-exchanges:limits}, the order rule of 100 per 10 seconds allows ten messages a second, but the daily rule of 200\,000 allows
only 2.3 a second sustained: a firm quoting five instruments can change each quote about once every two seconds on average, unless it negotiates
higher limits. Venues also publish several public endpoints of different characters.

\begin{dated}{2026-09}{A venue's public endpoints}\label{dat:m3:getting-access-crypto-venues:endpoints}
Binance's spot API documentation lists base endpoints api.binance.com, api-gcp.binance.com and api1 to api4.binance.com, noting that the last four
``should give better performance but have less stability'', a separate endpoint for public market data only, and a request validity window
(\texttt{recvWindow}) of 5\,000 milliseconds by default and 60\,000 at most.
\end{dated}

\section{Loan lines, off-exchange settlement and custody}

\begin{definition}[Loan line, collateral mirroring]\label{def:m3:getting-access-crypto-venues:loan}
A \emph{loan line}\index{loan line} is a credit facility from a venue or its affiliate that lets an institutional client borrow assets or trade on margin
beyond its deposits, against collateral and a credit assessment. \emph{Collateral mirroring}\index{collateral mirroring} is the crediting, on a venue, of a
trading balance equal to collateral that the client keeps with a custodian, which is the mechanism of \omterm{def:m3:centralised-exchanges:oes}{off-exchange settlement}
(\cref{ch:m3:centralised-exchanges}).
\end{definition}

The access decision is therefore also a credit decision in both directions. The firm extends credit to the venue by depositing; the venue extends credit
to the firm through \omterm{def:m3:getting-access-crypto-venues:loan}{loan lines} and mirrored balances, and prices it in the limits and haircuts it applies. A firm that trades on five venues with
\omterm{def:m3:centralised-exchanges:oes}{off-exchange settlement} through one custodian holds most of its collateral in one place, nets its exposures across venues, and replaces five venue
exposures with one custodian exposure and the unsettled balances in between.

\section{Tutorial: five venues, one budget}\label{tut:m3:getting-access-crypto-venues:fees}

\begin{tutorial}
\textbf{Goal.} Given a month's volume and maker share across five venues and their fee schedules, compute each venue's fees, the saving from
concentrating volume, the value of a tier match, and the uptime a strategy achieves against a programme's conditions. \textbf{End state:}
\cref{fig:m3:getting-access-crypto-venues:fees,fig:m3:getting-access-crypto-venues:uptime} and the numbers of the weekend problem.
\begin{tutsteps}
\item \textbf{Tiers and fees.} The tier for a volume, a month's bill, and a matched tier.
\omcode{code/firm/mmprogram/firm_mmprogram.py}{18}{34}{Volume tiers, a month's fees and a tier match.}\label{lst:m3:getting-access-crypto-venues:fees}
\item \textbf{Uptime.} Sampled snapshots of the firm's own quotes checked against spread and depth.
\omcode{code/firm/mmprogram/firm_mmprogram.py}{46}{58}{Compliance of a snapshot, uptime and the programme rebate.}\label{lst:m3:getting-access-crypto-venues:uptime}
\item \textbf{Run} \texttt{spread\_vs\_concentrate()}, \texttt{tier\_match\_value()}, \texttt{programme\_value} and \texttt{fig\_cryptoaccess.py}.
\end{tutsteps}
\textbf{What to change next.} Find the allocation over the five venues that minimises fees subject to at least 10\% on each (for resilience); add a
programme's fixed fee; recompute uptime when the \omterm{def:m3:centralised-exchanges:ratelimit}{rate limit}, not volatility, is what makes quotes stale.
\end{tutorial}

\section{Build: the market-maker programme monitor}\label{bld:m3:getting-access-crypto-venues:mmprogram}

\begin{build}
\bfield{Purpose} The miniature firm must know, every day, which tier it is on at each venue and what its allocation costs, whether it is meeting each
programme's obligations, and how many quote changes its limits allow.
\bfield{Interface} \texttt{VolTier(min\_volume, maker\_bp, taker\_bp)}; \texttt{tier\_for(volume, tiers)}; \texttt{monthly\_fees(volume, maker\_share,
tiers, matched)}; \texttt{Snapshot(mid, bid, ask, bid\_depth, ask\_depth)}; \texttt{compliant}; \texttt{uptime}; \texttt{rebate}; \texttt{requote\_budget(governor,
instruments)} from \texttt{firm.ratelimit}.
\bfield{Rules} Schedules versioned by date; volume counted as the venue counts it; snapshots of the firm's own quotes, sampled as the venue samples; a
programme's rebate booked only when its measurement confirms it.
\bfield{Acceptance tests} \texttt{code/firm/mmprogram/tests/}: tiers and a month's bill; a tier match that helps only when better; compliance, uptime and
the all-or-nothing rebate; the requote budget set by the tightest order rule.
\bfield{Stretch} An optimiser for the allocation of volume across venues; intraday uptime alerts; the programme's measurement reproduced from the venue's
market data.
\end{build}

\begin{omsources}
\item Binance, spot fee schedule and spot API documentation (general information, limits); Kraken, fee schedule. Accessed September 2026.
\item Copper, ClearLoop product page (see \cref{ch:m3:centralised-exchanges}).
\item Kraken, support article ``How to get verified for a Business account'' (updated 24 August 2026); Bitget, press release of 16 September 2026;
Binance, press release of 4 June 2025 (PR Newswire); Bybit, help centre, ``Introduction to the Market Maker Incentive Program'' and
``Bybit Trading Fee Structure'', September 2026.
\end{omsources}

\section{Exercises}

\begin{exercise}[$\star$]\label{exo:m3:getting-access-crypto-venues:1}
On Binance's schedule, what does a month of USD~100 million, 70\% as maker, cost?
\end{exercise}

\begin{exercise}[$\star$]\label{exo:m3:getting-access-crypto-venues:2}
What does a venue's \omterm{def:m3:getting-access-crypto-venues:kyc}{know-your-customer check} ask of a trading firm, and why does the venue care?
\end{exercise}

\begin{exercise}[$\star$]\label{exo:m3:getting-access-crypto-venues:3}
Why does the daily order rule, not the ten-second one, set a quoting firm's sustained budget at the limits of
\cref{dat:m3:centralised-exchanges:limits}?
\end{exercise}

\begin{exercise}[$\star\star$]\label{exo:m3:getting-access-crypto-venues:4}
Venue C matches the firm's Binance tier on its USD~120 million there. What is the match worth a month?
\end{exercise}

\begin{exercise}[$\star\star$]\label{exo:m3:getting-access-crypto-venues:5}
A programme pays 0.5 basis points on maker volume at 90\% uptime. On USD~210 million of maker volume, what is it worth, and what is the cost of
missing the requirement by one point?
\end{exercise}

\begin{exercise}[$\star\star$]\label{exo:m3:getting-access-crypto-venues:6}
Why might a firm keep 10\% of its volume on each of five venues even though concentrating it is cheaper?
\end{exercise}

\begin{exercise}[$\star\star\star$]\label{exo:m3:getting-access-crypto-venues:7}
\emph{Coding.} With \texttt{programme\_value}, find the uptime at maximum spreads of 10 and 20 basis points when depth is halved above 2.5 times normal
volatility.
\end{exercise}

\begin{exercise}[$\star\star\star$]\label{exo:m3:getting-access-crypto-venues:8}
\emph{Find the flaw.} ``We will reach the VIP tier faster by trading with ourselves across two \omterm{def:m3:centralised-exchanges:subaccount}{sub-accounts}.''
\end{exercise}

\section{Problem: Five Venues, One Budget}

\begin{problem}[{Weekend problem --- allocating a month of volume}]\label{pb:m3:getting-access-crypto-venues:1}
A firm trades USD~600 million a month, 70\% as maker. It can spread the volume evenly over Binance, Kraken and three other venues (C, D and E, on the
tutorial's illustrative schedules), or concentrate it on Binance and Kraken. One venue offers a programme paying 0.5 basis points on maker volume for
90\% uptime at a maximum spread of 20 basis points and USD~50\,000 a side.

\medskip\noindent\textbf{Part I --- The bill.}
\begin{enumerate}
\item What does each venue cost when the volume is spread evenly, and in total?
\item What do the two venues cost when the volume is concentrated?
\item What is the monthly saving?
\item Which venues' tiers change, and why is Kraken so much cheaper for a maker?
\item What would a tier match at venue C change?
\end{enumerate}

\medskip\noindent\textbf{Part II --- The programme.}
\begin{enumerate}[resume]
\item On USD~300 million at the programme's venue, what is the rebate worth?
\item What uptime does the firm's strategy achieve, and does it qualify?
\item What change to the strategy lifts it past 90\%?
\item What does that change cost in risk?
\item How does the requirement interact with the venue's \omterm{def:m3:centralised-exchanges:ratelimit}{rate limits}?
\end{enumerate}

\medskip\noindent\textbf{Part III --- Beyond fees.}
\begin{enumerate}[resume]
\item What exposure to the venues does each allocation create?
\item How would \omterm{def:m3:centralised-exchanges:oes}{off-exchange settlement} change the comparison?
\item What does the firm give up by leaving three venues?
\item How should it measure where its trades are best executed?
\item How often should it revisit the allocation?
\end{enumerate}

\medskip\noindent\textbf{Part IV --- Judgement.}
\begin{enumerate}[resume]
\item Is concentrating volume good for the market?
\item Why do venues pay rebates rather than cut fees for everyone?
\item What is the right number of venues for a new firm?
\item State the \emph{named result}: the monthly saving from concentrating on two venues, and the uptime the firm must keep, and currently keeps, for the rebate.
\item In one sentence: what does a crypto venue sell a market maker?
\end{enumerate}
\end{problem}

\section{Interview questions}

\begin{interviewq}[$\star$ \iqroles{trader}]\label{iq:m3:getting-access-crypto-venues:1}
How do volume tiers change where you should trade?
\end{interviewq}

\begin{interviewq}[$\star$ \iqroles{risk}]\label{iq:m3:getting-access-crypto-venues:2}
What would you check before depositing USD~20 million on a new venue?
\end{interviewq}

\begin{interviewq}[$\star\star$ \iqroles{developer}]\label{iq:m3:getting-access-crypto-venues:3}
How do you design a quoting system to meet an \omterm{def:m3:getting-access-crypto-venues:uptime}{uptime requirement} without exceeding \omterm{def:m3:centralised-exchanges:ratelimit}{rate limits}?
\end{interviewq}

\begin{interviewq}[$\star\star$ \iqroles{trader}]\label{iq:m3:getting-access-crypto-venues:4}
Is a market-maker rebate worth taking if it requires quoting through volatile periods?
\end{interviewq}

\begin{interviewq}[$\star\star$ \iqroles{researcher}]\label{iq:m3:getting-access-crypto-venues:5}
How would you measure whether a venue's volume is real before joining its programme?
\end{interviewq}

\begin{interviewq}[$\star\star\star$ \iqroles{developer, trader}]\label{iq:m3:getting-access-crypto-venues:6}
Build the daily report a head of trading needs on the firm's five venue relationships.
\end{interviewq}
